
\chapter{Programming Languages and Memory Semantics for TM}
\section{C/C++ Memory Model and ABI Hazards: \texttt{volatile}, Inline Assembly, and Atomics}
\begin{itemize}
    \item \texttt{volatile} only prevents compile‑time reordering, not CPU reordering.
    \item C11/C++11 atomics and memory ordering (relaxed, acquire, release, seq\_cst).
    \item How compiler barriers can replace costly full fences inside a TM.
    \item The per-operation cost model in STM backends: which ordering each read/write actually needs, and how backends downgrade \texttt{seq\_cst} to \texttt{compiler\_fence(SeqCst)} on the read path.
    \item Comparing the TM read path to lock-free readers/`-writer patterns: same atomics, different workload.
    \item ABI hazards specific to C/C++ TM integration: calling-convention mismatches between instrumented and non-instrumented functions, and the DATA-vs-TEXT symbol traps that arise when a runtime exposes hooks as function pointers rather than functions (see the pitfalls chapter).
\end{itemize}

\begin{figure}[htbp]
\centering
\resizebox{\textwidth}{!}{%
\begin{tikzpicture}[
    node distance=1.55cm and 1.85cm,
    box/.style={draw, rounded corners, align=center, minimum width=2.7cm, minimum height=0.85cm},
    small/.style={draw, rounded corners, align=center, minimum width=2.35cm, minimum height=0.75cm},
    arr/.style={-{Latex[length=2mm]}, thick}
]
\node[box] (src) {source\\\texttt{\_\_transaction\_atomic}};
\node[box, right=of src] (pass) {compiler pass\\address check};
\node[box, right=of pass] (abi) {ABI boundary\\real functions};
\node[box, right=of abi] (rt) {TM runtime\\read/write API};

\node[small, below=of pass] (plain) {plain load/store\\if private};
\node[small, below=of rt] (atoms) {atomics\\acquire/release};
\node[small, right=of atoms] (fence) {compiler fences\\read path};

\draw[arr] (src) -- (pass);
\draw[arr] (pass) -- node[above]{shared} (abi);
\draw[arr] (abi) -- (rt);
\draw[arr] (pass) -- node[left]{private} (plain);
\draw[arr] (rt) -- (atoms);
\draw[arr] (rt) -- (fence);
\draw[arr] (atoms) -- (fence);
\end{tikzpicture}%
}
\caption{A language-integrated TM boundary: the compiler decides which memory operations are instrumented; the runtime decides which atomics and fences implement the protocol.}
\label{fig:language-tm-boundary}
\end{figure}

\section{Language Integration Case Studies}
\subsection{Rust: The Borrow Checker and TM}
\begin{itemize}
    \item The borrow checker eliminates data races at compile time.
    \item \texttt{UnsafeCell} and interior mutability: what remains unprotected.
    \item How Rust’s safety guarantees can reduce (but not eliminate) the need for TM instrumentation.
    \item Rust's guarantees eliminate \emph{data races} but not \emph{deadlocks} or \emph{lost atomicity} --- the very guarantees TM adds on top of the language.
    \item \texttt{Sync}/\texttt{Send} traits as a way to reason about which values can be shared across threads; TM reifies ``shared mutable'' safety at runtime instead of compile time.
    \item The C++ and Rust stories are different in kind, not just degree: C++ gives the programmer \emph{control} over orderings; Rust removes whole classes of hazard at compile time. A TM layer must therefore integrate with each differently --- one C++ ABI, one Rust trait boundary.
\end{itemize}

\subsection{Language‑Integrated Annotations for TM}
\begin{itemize}
    \item GCC’s \texttt{\_\_transaction\_atomic} and \texttt{transaction\_safe} / \texttt{transaction\_pure} attributes.
    \item The (abandoned) C++ TM Technical Specification: transaction-safe functions, atomic blocks, and why the committee ultimately dropped it.
    \item Research proposals like \texttt{@atomic} blocks.
    \item This book's instrumentation plugin: a \texttt{[[tm::shared]]} annotation on globals, and the \texttt{LLVM\_TM\_ADDR\_CHECK} threshold deciding when instrumentation is actually needed.
    \item The three case studies (GCC, the C++ TS, this book's plugin) are a spectrum of how much the language commits to: the less the compiler enforces, the more the runtime must prove.
\end{itemize}

\begin{table}[htbp]
\centering
\begin{tabular}{p{0.21\textwidth}p{0.25\textwidth}p{0.25\textwidth}p{0.20\textwidth}}
\hline
Mechanism & Programmer contract & Runtime/compiler burden & Main hazard \\
\hline
\texttt{volatile} & Preserve visible accesses at compile time & Almost none & Not synchronization; CPU may reorder \\
C/C++ atomics & Choose ordering per operation & Implement the requested ordering & Correct but non-composable protocols \\
GCC atomic block & Mark transactional regions and safe calls & Compiler lowers loads, stores, and calls & Unsafe escape through non-transactional code \\
\texttt{[[tm::shared]]} plugin & Mark shared objects explicitly & Instrument only addresses that may alias shared state & Missed annotation loses atomicity \\
Rust trait boundary & Share only values satisfying \texttt{Send}/\texttt{Sync} contracts & Encapsulate mutation behind safe APIs or \texttt{UnsafeCell} & Interior mutability and FFI escape hatches \\
\hline
\end{tabular}
\caption{Language mechanisms differ in where they place the TM proof obligation.}
\label{tab:language-tm-taxonomy}
\end{table}

\section{Worked Example: Atomic Operations for a Counter}
\index{atomic operations}\index{C11 atomics}\index{volatile}\index{atomic\_signal\_fence}
\label{sec:atomics-example}

A shared counter is the smallest TM-relevant program. Compare the three
synchronisation styles on the same three lines of logic:

\begin{lstlisting}[language=C, caption={A shared counter three ways.}]
/* Locks: correct, coarse, non-composable. */
pthread_mutex_lock(&m); counter++; pthread_mutex_unlock(&m);

/* C11 atomics: fine-grained, but the orderings are on the programmer. */
int prev = atomic_fetch_add(&counter, 1, memory_order_relaxed);

/* TM: the block is the unit; orderings are the runtime's problem. */
__transaction_atomic { counter++; }
\end{lstlisting}

\begin{itemize}
    \item The atomic version exposes \emph{which} ordering to use per operation; the TM version hides the whole protocol.
    \item The TM runtime still needs the right atomics internally --- it just puts them behind one API (Part~IV).
    \item Composability reappears: two counters in one \texttt{atomic} block compose; two separate \texttt{atomic\_fetch\_add} calls do not.
    \item Exercise~3 of this chapter asks for the full C++ and Rust comparison.
\end{itemize}

\section{Worked Example: ABI-Safe Bank Transfer}
\index{ABI hazards}\index{compiler fence}\index{single global lock}\index{bank transfer}
\label{sec:abi-safe-bank-transfer}

The following example is deliberately an SGL-style TM. Its purpose is not
performance; it is the language boundary. The runtime exports functions, not
data symbols containing function pointers. The instrumented code calls those
functions and lets the runtime choose atomics and fences internally.

\begin{lstlisting}[language=C++, caption={A compilable C++ TM facade for the bank-transfer invariant.}]
#include <array>
#include <atomic>
#include <cassert>
#include <cstddef>

struct Account {
    int balance;
};

static std::array<Account, 2> bank{{Account{100}, Account{100}}};
static std::atomic_flag global_lock = ATOMIC_FLAG_INIT;

/* ABI boundary: real functions, not DATA symbols containing function pointers. */
extern "C" void tm_begin() {
    while (global_lock.test_and_set(std::memory_order_acquire)) {
        /* spin */
    }
}

extern "C" void tm_commit() {
    global_lock.clear(std::memory_order_release);
}

extern "C" int tm_read_i32(const int *p) {
    std::atomic_signal_fence(std::memory_order_seq_cst);
    int v = *p;
    std::atomic_signal_fence(std::memory_order_seq_cst);
    return v;
}

extern "C" void tm_write_i32(int *p, int v) {
    std::atomic_signal_fence(std::memory_order_seq_cst);
    *p = v;
    std::atomic_signal_fence(std::memory_order_seq_cst);
}

class TxGuard {
public:
    TxGuard() { tm_begin(); }
    ~TxGuard() { tm_commit(); }

    TxGuard(const TxGuard&) = delete;
    TxGuard& operator=(const TxGuard&) = delete;
};

static bool transfer(std::size_t from, std::size_t to, int amount, int floor) {
    TxGuard tx;

    int from_balance = tm_read_i32(&bank[from].balance);
    int to_balance = tm_read_i32(&bank[to].balance);

    if (from_balance - amount < floor) {
        return false;
    }

    tm_write_i32(&bank[from].balance, from_balance - amount);
    tm_write_i32(&bank[to].balance, to_balance + amount);
    return true;
}

static int total_money() {
    return bank[0].balance + bank[1].balance;
}

int main() {
    const int initial_total = total_money();

    assert(transfer(0, 1, 30, 0));
    assert(bank[0].balance >= 0);
    assert(bank[1].balance >= 0);
    assert(total_money() == initial_total);

    assert(!transfer(0, 1, 1000, 0));
    assert(total_money() == initial_total);
}
\end{lstlisting}

\begin{itemize}
    \item The acquire on \texttt{tm\_begin} and release on \texttt{tm\_commit} publish the critical section; the per-access fences only constrain the compiler.
    \item The source-level transaction is compositional: the floor check and both writes are one unit.
    \item The ABI contract is visible: instrumented code calls \texttt{tm\_read\_i32} and \texttt{tm\_write\_i32}; non-instrumented code must not bypass shared state.
    \item This is the semantic baseline implemented by SGL and TSX-SGL before more selective protocols such as NOrec, TL2, and TSC-TM reduce contention.
\end{itemize}

\section{Worked Example: Model Checking Lost Atomicity}
\index{TLA+}\index{lost atomicity}\index{money conservation}
\label{sec:tla-bank-language-contract}

Language support for TM should preserve the same invariant as the runtime:
the transfer either commits as one state transition or has no effect. A small
TLA+ model captures that contract without mentioning cache lines, fences, or
calling conventions.

\begin{lstlisting}[language={}, caption={A TLC-checkable TLA+ model of an atomic bank transfer.}]
---- MODULE BankTMAtomicity ----
EXTENDS Naturals

A0 == 10
A1 == 10
Floor == 0
Amount == 3

VARIABLES bal, pc

vars == <<bal, pc>>

Init ==
    /\ bal = [0 |-> A0, 1 |-> A1]
    /\ pc = "idle"

CanTransfer ==
    bal[0] >= Amount + Floor

Begin ==
    /\ pc = "idle"
    /\ pc' = "active"
    /\ UNCHANGED bal

Commit ==
    /\ pc = "active"
    /\ CanTransfer
    /\ bal' = [bal EXCEPT ![0] = @ - Amount,
                          ![1] = @ + Amount]
    /\ pc' = "idle"

Abort ==
    /\ pc = "active"
    /\ ~CanTransfer
    /\ pc' = "idle"
    /\ UNCHANGED bal

Idle ==
    /\ pc = "idle"
    /\ UNCHANGED vars

Next ==
    Begin \/ Commit \/ Abort \/ Idle

Spec ==
    Init /\ [][Next]_vars

InvMoney ==
    bal[0] + bal[1] = A0 + A1

InvFloor ==
    bal[0] >= Floor /\ bal[1] >= Floor

THEOREM Spec => []InvMoney
THEOREM Spec => []InvFloor
====
\end{lstlisting}

\begin{itemize}
    \item The model treats the transaction body as one commit transition; this is the property language integration must preserve.
    \item \texttt{InvMoney} catches partial instrumentation: debiting without crediting violates conservation.
    \item \texttt{InvFloor} catches a misplaced check: the account floor must be tested in the same atomic unit as the debit.
    \item The model abstracts away C++ and Rust details; those details re-enter at the ABI and trait boundaries.
\end{itemize}

\section{Summary}
\begin{itemize}
    \item Different languages offer different levels of memory‑ordering control.
    \item A type system that prevents data races still cannot enforce atomicity across multiple memory locations without TM support.
\end{itemize}

\section{Exercises}
\begin{enumerate}
    \item In C, a compiler hoists a read of \texttt{x} out of an \texttt{atomic} block. Explain how \texttt{volatile} or an \texttt{asm} barrier can prevent this, and what the runtime cost is.
    \item A Rust program uses \texttt{RefCell}. Explain why an STM that instruments all loads and stores still cannot protect against a data race without additional runtime checks.
    \item Compare the programming effort required to add a thread-safe counter using locks, atomics, and a TM in C++ and Rust.
    \item \emph{[implementation]} Extend the ABI-safe bank-transfer example with a second shared object, \texttt{transfer\_count}. Update it inside the same transaction as the two account balances. Then write a non-transactional version using three separate atomics and identify an interleaving that preserves each individual atomic operation but violates the intended compound invariant.
    \item \emph{[verification]} Modify the TLA+ model in the lost-atomicity example by splitting \texttt{Commit} into two transitions, one debit and one credit. Run TLC on the finite model and record the shortest counterexample to \texttt{InvMoney}. Explain which source-level compiler or ABI mistake the split transition represents.
    \item \emph{[theory]} Compare a C++ TM API based on explicit \texttt{std::memory\_order} parameters with a Rust API based on a \texttt{TxCell<T>} wrapper. For each API, state what the type system proves, what the runtime must still validate, and where unsafe escape hatches can break atomicity.
\end{enumerate}

%  PART III: FORMAL FOUNDATIONS: MATHEMATICS AND VERIFICATION
